\section{Simple bundles with a prescribed middle Chern class}
\label{sec:simple-hodge-bundles}

This section constructs simple bundles on a CM product whose entire
Chern character stays of Hodge type on a polarized Weil family.
The construction solves the existence problem for the bundle at that
product point.  Its deformation is a further condition, and the
particular gluing used below fails that condition.

For $n\ge2$ put $g=2n$ and define
\begin{equation}
 F_n(t)=\sum_{k=0}^{2n}
  \frac{e_k(1,-1,\ldots,n,-n)}{(2n)_k}t^k,
 \qquad (2n)_k=\frac{(2n)!}{(2n-k)!}.
 \label{eq:positive-character-polynomial}
\end{equation}
Here $e_k$ is the elementary symmetric polynomial and $(2n)_0=1$.

\begin{theorem}\label{thm:simple-hodge-bundle}
Let $K$ be any imaginary quadratic field, let $E$ be a complex
elliptic curve with $\operatorname{End}^0(E)=K$, and give
$A=E^{2n}$ its product principal polarization $\theta$.
For every nonzero rational primitive class
$w\in H^{n,n}(A)\cap H^{2n}(A,\mathbf Q)$ there are positive integers
$R,N$, a nonzero rational number $c$, and a simple vector bundle
$\mathcal E$ such that
\[
 \operatorname{ch}(\mathcal E)=R F_n(N\theta)+cw.
\]
It is an ordinary extension of two finite direct sums of line
bundles of $\theta$-slope zero.  The ratio $c/(RN^n)$ can be chosen
arbitrarily small and positive.  No bound for $R,N$ uniform in
$E,w$ is asserted.

If $w$ is a Weil class for a compatible action of signature $(n,n)$,
this entire character stays of Hodge type on the corresponding
polarized Weil domain.  For $n\ge3$ its degree-two Hochschild action
has rank $7n^2-2n$.  Equality
$\dim\operatorname{Ext}^2(\mathcal E,\mathcal E)=7n^2-2n$ would
therefore imply strong semiregularity; that equality is not part
of the construction.
\end{theorem}

The proof uses exact finite averages to make a signed cohomology
identity into a sum with positive multiplicities.  Rank-two
differences between the divisor classes then provide extension
groups on abelian surfaces.  Enough independent extension classes
force scalar endomorphisms of the resulting bundle.
The general Hochschild-rank formula for a Weil class plus a
polarization polynomial is
\cite[Theorem~A6]{BBK3Sixfolds}; we compute the relevant minor for
$F_n$ and give the resulting kernel calculation explicitly.

\subsection{A rational unitary orbit and its primitive summand}

Use the central CM action on the product to identify
$\operatorname{NS}(A)_{\mathbf Q}$ with the Hermitian matrices
$\operatorname{Herm}_g(K)$.  This follows directly from the
Rosati-fixed endomorphisms of the product polarization.
This central action is used to construct divisors; it need not be
the balanced Weil action in the theorem.  The class $\theta$
corresponds to the identity matrix.  Put
\[
 D=\operatorname{diag}(1,-1,\ldots,n,-n).
\]
Every $D_u=uDu^*$ for $u\in U(g,K)$ is a rational divisor class
of trace, and hence slope, zero.  This rational unitary group is
dense in $U(g)$: skew-Hermitian matrices over $K$ are dense in all
skew-Hermitian complex matrices, and the Cayley transform covers
a dense chart of the unitary group.

We need an assertion about the whole character $e^{D_u}$, not just
its middle component.  As a complex unitary representation,
\[
 H^{p,p}(A)=\Lambda^p\mathbf C^g\otimes
                         (\Lambda^p\mathbf C^g)^*.
\]
Its irreducible summands have highest weights
$(1^k,0^{g-2k},-1^k)$ for $0\le k\le\min(p,g-p)$, each once.
Indeed write the dual factor as
$\Lambda^{g-p}\mathbf C^g\otimes\det^{-1}$.
The exterior Pieri rule gives the two-column diagrams
$(2^k,1^{g-2k})$, and subtracting the determinant weight gives
the displayed list.  The summand with $k=n$ occurs only when
$p=n$.  It is exactly
\[
 P^{n,n}=\ker(\theta\wedge:H^{n,n}\longrightarrow H^{n+1,n+1}).
\]
It is irreducible and occurs once in the entire graded Hodge space.

Its projection from $e^D$ is nonzero.  Write $x_1,\ldots,x_{2n}$
for the coordinate elliptic divisor classes, so that $x_j^2=0$ and
$\int_Ax_1\cdots x_{2n}=1$.  The rational class
\[
 q=\prod_{j=1}^n(x_{2j-1}-x_{2j})
\]
is primitive, and
\[
 \int_A\frac{D^n}{n!}q=(-2)^n n!\ne0.
\]
Complete reducibility and the unique occurrence of $P^{n,n}$ imply
that the real span of the orbit $\{e^{D_u}:u\in U(g)\}$ contains
every primitive middle class with all its other graded components
zero.  To see why multiplicity matters, the span is a submodule;
its nonzero projection to an irreducible constituent occurring once
contains that constituent itself.

The rational span of the rational-unitary orbit has closed real
span, hence contains the continuous orbit by density.  All rational
Hodge classes on a CM elliptic product have the indicated rational
form: the diagonal CM Hodge torus fixes exactly the equal
holomorphic and antiholomorphic degrees.  Rational row reduction
therefore gives a finite exact equality
\begin{equation}
 w=\sum_a b_a e^{D_a},\qquad b_a\in\mathbf Q,
 \quad D_a=u_aDu_a^*,\quad u_a\in U(g,K).
 \label{eq:primitive-orbit-relation}
\end{equation}
In particular no approximation to $w$ is involved.

\subsection{Finite averages make all multiplicities positive}

In each unitary eigenbasis permute the distinct eigenvalues of $D$.
Let $\mathcal O_a$ be the resulting set of $g!$ divisor classes.
The rank-one divisor forms $\omega_1,\ldots,\omega_g$ of that basis
satisfy $\omega_j^2=0$ and $\sum_j\omega_j=\theta$.
Counting the subsets in the exponential expansion gives
\begin{equation}
 \frac1{g!}\sum_{D'\in\mathcal O_a}e^{D'}
 =\sum_{k=0}^g\frac{e_k(1,-1,\ldots,n,-n)}{\binom gk}
                              \frac{\theta^k}{k!}
 =F_n(\theta).
 \label{eq:finite-positive-average}
\end{equation}
If a coefficient in \eqref{eq:primitive-orbit-relation} is negative,
add its absolute value times the full corresponding orbit sum to
both sides.  This removes that negative coefficient and introduces
no new negative ones.  There are finitely many coefficients.
Then add a strictly positive constant on every orbit involved.
After identifying repeated classes the result is
\begin{equation}
 C F_n(\theta)+w=\sum_v q_v e^{D_v},
 \qquad C>0,\quad q_v\in\mathbf Q_{>0}.
 \label{eq:positive-orbit-relation}
\end{equation}
Additional uniform orbit sums make $1/C$ arbitrarily small.
The equality holds in every character degree, including degree zero.

We may connect these orbit sets while preserving
\eqref{eq:positive-orbit-relation}.  Inside an orbit, every coordinate
transposition gives a difference
$(\lambda_j-\lambda_k)(\omega_k-\omega_j)$ of Hermitian rank two
and signature $(1,1)$.  Its graph is the bipartite transposition
graph of the symmetric group.

Any rational unitary transformation is a product of rank-one unitary
transformations.  For completeness, if $v\in K^m$ is a unit vector
different from $e_1$, put $u=v-e_1$ and
\[
 h=I-\frac{uu^*}{1-v_1}.
\]
The identities $u^*v=1-v_1$ and
$\|u\|^2=2-v_1-\overline v_1$ give $hv=e_1$ and $hh^*=I$.
The denominator is nonzero, since $v_1=1$ for a unit vector forces
$v=e_1$.  Apply this operation to the first column, then induct on
its orthogonal complement.  For a rank-one unitary transformation,
\[
 hD'h^*-D'=(h-I)D'h^*+D'(h^*-I)
\]
has rank at most two and trace zero.  If nonzero, it has rank two
and signature $(1,1)$.

Include the finitely many intermediate eigenbasis orbits along
these products.  Give each a positive uniform coefficient using
\eqref{eq:finite-positive-average}.  Equal orbits can be identified:
one common matrix with distinct eigenvalues determines the same
unordered eigenlines and therefore the same permutation orbit.
Choose a tree joining the distinct orbit sets, and one of the
rank-two edges for each tree edge.  Keep all internal transposition
edges.  Flip the parity colorings successively along the tree.
The result is a finite connected bipartite graph of distinct rational
divisor classes, every edge having a difference of signature $(1,1)$
on a rational surface quotient.

\subsection{Line bundles, surface cohomology and simplicity}

Clear the denominators of the $q_v$ to obtain positive multiplicities
$m_v$.  Choose $N$ so that every alternating Riemann form of $ND_v$
has even integral values on the product lattice $\Lambda$.
For each choose its Appell--Humbert line bundle with trivial
semicharacter.  The cocycle formula and tensor additivity are
\cite[\S2, pp.~20--21]{MumfordAV}.
Explicitly the factor for a Hermitian form $H$ is
\[
 j_\lambda(z)=
 \exp\bigl(\pi H(z,\lambda)+\tfrac\pi2H(\lambda,\lambda)\bigr).
\]
Even integrality makes the omitted semicharacter equal to one
admissible.  This formula is additive in $H$, so the quotient of
two chosen line bundles has exactly the difference form, with
no extra class in $\operatorname{Pic}^0$.

For an edge let $W$ be the kernel of that difference.  It is rational
over $K$, so $B=W/(W\cap\Lambda)$ is an abelian subvariety and
$p:A\to Y=A/B$ is a quotient onto an abelian surface.
The form vanishes when either argument is in $W$; the displayed
cocycle therefore descends to $Y$.  Thus
\[
 L_vL_u^{-1}=p^*M
\]
for a nondegenerate line bundle of index one on $Y$.
Mumford's vanishing and index theorems
\cite[\S16, pp.~150,155,162]{MumfordAV} give
$H^i(Y,M)=0$ for $i\ne1$ and $h^1(Y,M)=-\chi(M)>0$;
the inverse has the same index.
The low-degree Leray sequence, $p_*\mathcal O_A=\mathcal O_Y$,
and $H^0(Y,M)=0$ give
\[
 H^1(A,p^*M)=H^1(Y,M).
\]
Here the higher direct images of $\mathcal O_A$ are trivial with
fibres $H^q(B,\mathcal O_B)$: invariant forms on $A$ restricting to
a basis on $B$ give a global frame.
Surface Serre duality also shows that the two reverse extension
groups pair into the nonzero line
\begin{equation}
 q_{uv}=p^*H^2(Y,\mathcal O_Y)\subset H^2(A,\mathcal O_A).
 \label{eq:edge-surface-line}
\end{equation}
Replacing $N$ by a positive multiple tensors $M$ by that power.
Surface Riemann--Roch makes its $h^1$ grow quadratically.
We can therefore arrange
\[
 \dim\operatorname{Ext}^1(L_u,L_v)\ge m_um_v
\]
on every edge in either direction.

Distinct $L_u,L_v$ have no homomorphisms in either direction.
Their divisor difference is nonzero of $\theta$-degree zero,
whereas every nonzero effective divisor has positive $\theta$-degree.
If a section had no divisor it would trivialize the quotient,
also contradicting the distinct Chern classes.

Let $P$ and $Q$ be the line-bundle sums with multiplicities $m_v$
on the two parts of the bipartition.  On every chosen edge from
$L_u\subset P$ to $L_v\subset Q$, choose $m_um_v$ independent
extension classes $\eta_{ab}$ and set
\[
 e_{vu}=\sum_{a,b}\eta_{ab}\otimes E_{ab},
\]
where $E_{ab}$ are the matrix units.  Set all other components zero.
This is a single class in $\operatorname{Ext}^1(P,Q)$, defining
an ordinary extension
\begin{equation}
 0\longrightarrow Q\longrightarrow\mathcal E
                   \longrightarrow P\longrightarrow0.
 \label{eq:positive-bundle-extension}
\end{equation}
An extension of vector bundles is locally free.  Every endomorphism
preserves $Q$ because $\operatorname{Hom}(Q,P)=0$, and is determined
by its pair on $Q,P$ because $\operatorname{Hom}(P,Q)=0$.
Those endomorphisms are block diagonal.  The compatibility equation
at an edge gives $b_vE_{ab}=E_{ab}t_u$ for every matrix unit, forcing
both endpoint matrices to be the same scalar.  Connectivity forces
one scalar everywhere.  Thus $\operatorname{End}(\mathcal E)=\mathbf C$.

Additivity of the Chern character in
\eqref{eq:positive-bundle-extension}, followed by degreewise scaling
in \eqref{eq:positive-orbit-relation}, gives the asserted character.
Its rank is $R=\sum_vm_v$ and its middle coefficient is a nonzero
rational multiple of $N^n$.  This proves the existence assertion.
It also proves slope semistability, since the two sums have slope
zero.  Stability is not asserted.

\subsection{The Hochschild kernel and the number 104}

Now suppose $w$ is a compatible rational Weil class and $n\ge3$.
Write $F_n(t)=1+p_2t^2+p_4t^4+\cdots$.  Direct expansion of
$\prod_{j=1}^n(1-j^2t^2)$ gives
\[
 p_2=-\frac{(n+1)(2n+1)}{12(2n-1)},\qquad
 p_4=\frac{(n+1)(2n+1)(5n+6)}{1440(2n-3)},
\]
and
\begin{equation}
 6p_4-p_2^2=
 \frac{(n+1)(2n+1)(40n^3+12n^2-22n+33)}
 {720(2n-3)(2n-1)^2}>0.
 \label{eq:positive-rank-minor}
\end{equation}

Under HKR write a Hochschild class of degree two as
$(\xi,\eta,\zeta)$ in
$H^0(\Lambda^2T_A)\oplus H^1(T_A)\oplus H^2(\mathcal O_A)$.
Put $b=\xi\lrcorner\theta^2$ and $u=\eta\lrcorner\theta$.
The first three equations for annihilating $F_n(N\theta)$ are
\[
 p_2N^2b+\zeta=0,\qquad 2p_2N^2\theta u=0,
 \qquad 6p_4N^4\theta^2b+p_2N^2\theta^2\zeta=0.
\]
The identities used are the derivation identity for $\eta$ and
$\xi\lrcorner\theta^k=\binom k2\theta^{k-2}b$.
Lefschetz injectivity and \eqref{eq:positive-rank-minor} give
$b=u=\zeta=0$.  Contraction with $\theta^2$ identifies bivectors
with $H^{0,2}$, so $\xi=0$.  Conversely $u=0$ kills every power
of $\theta$.  Thus the kernel for the background polynomial is
exactly the polarized tangent space, of dimension $2n^2+n$.

The norm-one torus of the balanced $K$ action fixes $\theta$.
Its weights on $HH^2$ are $-2,0,2$, and its weights on the two Weil
lines are $2n,-2n$.  The polynomial output weights and the Weil
output weights are disjoint for $n\ge3$.  There is consequently
no cancellation between them, and the kernel for
$R F_n(N\theta)+cw$ is the intersection of their two kernels.

Split $H^{1,0}=U\oplus V$ according to the balanced action, with
$\dim U=\dim V=n$.  The polarized tangent space is
\[
 \operatorname{Sym}^2U^*\oplus(U^*\otimes V^*)
                         \oplus\operatorname{Sym}^2V^*.
\]
The middle summand kills both Weil determinant forms.  The other
two summands act injectively on their respective determinant forms:
contracting a holomorphic volume form identifies its next exterior
power with the dual standard representation and records the
symmetric matrix.  The outputs have distinct $K$ weights.
Both determinant coefficients of a nonzero rational $w$ are nonzero.
The common kernel is therefore exactly $U^*\otimes V^*$,
of dimension $n^2$.  Since $\dim HH^2=8n^2-2n$, the rank is
$7n^2-2n$ as asserted.

For an eightfold this becomes
\[
 F_4(t)=1-\frac{15}{28}t^2+\frac{91}{560}t^4
               -\frac{41}{1008}t^6+\frac1{70}t^8,
 \qquad 7\cdot4^2-2\cdot4=104.
\]
The trace/evaluation factorization of Buchweitz--Flenner
\cite[Proposition~6.2.1 and Theorem~6.4.2]{BFHH} identifies the
composite through $\operatorname{Ext}^2$ with this Hochschild action.
If that Ext group has dimension equal to the computed rank,
the semiregularity map is injective.  This proves the remaining
assertions of Theorem~\ref{thm:simple-hodge-bundle}.

\subsection{The fourfold case and its Euler obstruction}

The construction also works for $n=2$, but the weight separation
just used does not.  Here
$F_2(t)=1-5t^2/12+t^4/6$.
For a sufficiently small nonzero rational $\varepsilon$, the action
on $F_2(\theta)+\varepsilon w$ has rank $24$.
To check this, the $H^2(\mathcal O)\oplus H^0(\Lambda^2T)$ block
at $\varepsilon=0$, after Lefschetz identifications, has matrix
\[
 \begin{pmatrix}1&2p_2\\p_2&12p_4\end{pmatrix}\otimes I_6,
 \qquad 12p_4-2p_2^2=119/72.
\]
It remains invertible for small $\varepsilon$.
The $H^1(T)$ block has rank six at zero, with image
$\theta H^{0,2}$ and kernel the ten-dimensional polarized tangent.
On that kernel the Weil derivative has rank six, with kernel the
four-dimensional Weil tangent.  Its image is primitive, since
both $\kappa\theta=0$ and $\theta w=0$.
It is therefore disjoint from $\theta H^{0,2}$.
A twelve-by-twelve minor of the perturbed block has nonzero leading
term of order $\varepsilon^6$, so this block has rank twelve for
all sufficiently small nonzero $\varepsilon$.  The four-dimensional
Weil tangent bounds its rank above by twelve.  The total is $24$.
The positive completion permits such $\varepsilon$; subsequent
integral scaling conjugates these maps by invertible degree scalings.

None of these simple fourfold bundles is strongly semiregular.
On a Calabi--Yau fourfold the full degree-two trace map is dual to
degree-two evaluation.  If it were injective, evaluation would be
surjective, and their composite would have rank
$\operatorname{ext}^2=24$.
The degree-one action on a nonzero Weil class is injective; its
weights $\pm3$ are disjoint from the polynomial output weights
$\pm1$.  Thus $\operatorname{ext}^1\ge8$.
Simplicity and Serre duality would give
$\chi(\mathcal E,\mathcal E)=2-2\operatorname{ext}^1+
\operatorname{ext}^2\le10$.
But Riemann--Roch, $\theta w=0$, and Hodge--Riemann give
\[
 \chi(\mathcal E,\mathcal E)
 =\frac{73}{6}R^2N^4+c^2\int_Aw^2>10.
\]
This contradiction concerns ordinary Ext groups.  It makes no
assertion about invariants under an additional finite group.
In dimension eight the middle Ext group prevents the same Euler
argument from giving a contradiction.
